% Section 1: problem statement and motivation.
\section{Problem statement and motivation}\label{sec:intro}

\subsection{The problem statement}
MRPL's problem statement SIH26119, \emph{Indigenous GPU-Accelerated Optimization Solver (Sovereign
Alternative to Xpress/CPLEX)}, starts from a practical dependency. Refinery crude selection,
blending, production planning, scheduling and product distribution in India are formulated as
linear and mixed-integer programmes and solved with a small number of foreign commercial solvers.
Those solvers are closed: their licences cost money every year, their algorithms cannot be
inspected or adapted, and their availability depends on vendors outside the country. The
statement asks for a solver that is developed here, handles the same problem classes, can use
GPUs, and does not build on any existing solver library.

We read four requirements out of it:
\begin{enumerate}[label=\textbf{R\arabic*}]
  \item \textbf{Problem classes:} linear programmes (LP), mixed-integer linear programmes (MILP)
        and convex quadratic programmes (QP), the classes that planning and scheduling models
        fall into.
  \item \textbf{Independence:} no dependence on an existing optimisation solver or sparse direct
        solver library. The factorisations, orderings and algorithms must be our own code.
  \item \textbf{Interoperability:} read the MPS and QPS files that modelling tools and every
        public benchmark library already produce.
  \item \textbf{Acceleration:} use a GPU where the algorithm can exploit one.
\end{enumerate}

\subsection{What \nirnay{} is}
\nirnay{} (Sanskrit/Hindi \emph{nirṇay}, ``decision'') is a Python package whose compute
kernels are compiled to machine code with Numba and, for the GPU path, to CUDA with CuPy. It has
three LP engines, which suit different problem shapes:
\begin{itemize}
  \item a \textbf{bounded dual simplex} (\cref{sec:simplex}), which produces vertex solutions and
        re-solves quickly after a bound change: this is the engine branch-and-bound uses;
  \item a \textbf{primal--dual interior-point method} (\cref{sec:ipm}), whose iteration count
        grows only slowly with problem size;
  \item \textbf{PDLP}, a restarted primal--dual hybrid gradient method (\cref{sec:pdlp}). It needs
        only products with $A$ and $A\T$, so it runs on a GPU and scales to LPs too large to
        factorise.
\end{itemize}
On top of these sit a \textbf{branch-and-bound} MILP solver (\cref{sec:mip}) and a \textbf{QP
interior-point} method (\cref{sec:qp}). \Cref{fig:arch} shows how the parts fit together.

\subsection{Design principles}
\begin{description}[style=nextline,leftmargin=1.2em,font=\sffamily\bfseries\color{navy}]
  \item[Implement, do not wrap.] Every factorisation (LU, Cholesky, $LDL\T$), ordering and
        solver loop is in the repository. The code depends on NumPy arrays and the Numba
        compiler, plus CuPy for the GPU path. It depends on no optimisation or sparse-direct
        library. \Cref{tab:modules} lists the modules with their sizes.
  \item[Follow the primary literature.] Each module's docstring cites the papers its algorithm
        comes from, and this report follows those citations. Where the code departs from a
        paper, the report says so.
  \item[Verify every answer independently.] Benchmarks run every instance in its own process
        and re-solve it with a reference solver (HiGHS for LP and MILP, Clarabel for QP) that
        reads the file with its own parser (\cref{sec:verification}).
  \item[Report failures.] Tables in this report list every instance, solved or not, and
        \cref{sec:results} analyses the failures.
\end{description}

\subsection{Contributions and outline}
\Cref{sec:formulation} fixes the mathematical forms. \Cref{sec:architecture} describes the system
and the path of a solve, and \cref{sec:presolve} the presolve and postsolve. \Cref{sec:linalg}
covers the sparse linear algebra on which every engine rests. \Cref{sec:simplex,sec:ipm,sec:pdlp}
describe the three LP engines, \cref{sec:mip} branch-and-bound and \cref{sec:qp} the QP method.
\Cref{sec:safeguards} collects the numerical failures we met and how each was fixed.
\Cref{sec:verification} explains how answers are checked, \cref{sec:results} gives the benchmark
results, \cref{sec:cases} the industrial case studies built from Indian public data,
\cref{sec:competitors} a hands-on comparison with other solvers, and \cref{sec:roadmap} the
limitations and next steps. The appendices list every benchmark instance and the typesetting
toolchain of this report.
